\section{Market Commentary and Outlook}

\subsection{July in Review}

Digital asset markets extended their second-quarter recovery through July, with the Nexa Composite Digital Asset Index posting its fourth consecutive positive month despite a mid-month bout of volatility tied to broader risk-asset repricing. On-chain activity metrics -- active addresses, settled transaction value, and layer-2 sequencer throughput -- all reached multi-quarter highs during the month, which the Investment Committee interprets as evidence of genuine usage growth rather than purely speculative positioning. Derivatives markets remained orderly throughout the period: perpetual funding rates stayed within a normal band for the majority of the month, and open interest growth was accompanied by a proportionate rise in spot volume rather than the funding-rate spikes that typically precede sharp deleveraging events.

The Fund's overweight to layer-2 scaling infrastructure was the largest single contributor to July performance. Elevated sequencer throughput ahead of a widely anticipated network upgrade drove renewed attention to the segment, and the basket's constituents outperformed the broader benchmark by a wide margin over the month. The Investment Committee had built this position incrementally over the prior two quarters, sizing it according to its conviction-weighted framework as on-chain fundamentals strengthened.

\subsection{Positioning and Outlook}

Entering the third quarter, the Investment Committee maintains a moderately constructive stance, reflected in gross exposure of 96.8\% of NAV against a mandate ceiling of 130\%. Three considerations inform this positioning:

\begin{itemize}[leftmargin=1.4em,itemsep=2pt,topsep=2pt]
\item \textbf{Macro liquidity conditions.} Broad financial conditions have remained supportive through the third quarter to date, though the Committee continues to monitor central bank commentary for signs of a hawkish pivot that has historically preceded digital asset drawdowns.
\item \textbf{Network fundamentals.} Continued strength in on-chain settlement volume and layer-2 adoption supports the Fund's tactical overweight to scaling infrastructure, though the Committee will reassess sizing following the completion of the anticipated network upgrade.
\item \textbf{Valuation and positioning.} Derivatives positioning data show funding rates and open interest at levels the Committee characterizes as constructive rather than euphoric, leaving room for continued participation without the excess leverage that has historically amplified drawdowns.
\end{itemize}

The Committee intends to maintain the liquidity sleeve within its target range of 5\% to 12\% of NAV through the third quarter, preserving capacity for opportunistic rebalancing should realized volatility increase materially from current levels. No structural changes to the Fund's sleeve framework are planned for the coming quarter.
